\section{Core Theory and Concepts}
\label{sec:core-theory}

This section summarizes the underlying concepts of the technologies used to build the pipeline in Chapter~\ref{ch:proposed-solution}.

\subsection{Why Real-Time Processing Matters}

A modern platform such as a music-streaming service does not produce data in daily batches: it produces a continuous stream of small events -- a song played, a page viewed, a login -- the instant a user acts. Streamify, the project this report builds, simulates exactly this kind of platform with Eventsim, which generates synthetic ``song played'', ``page view'', and ``authentication'' events for a fake music website at a steady, ongoing rate.

In a traditional batch-processing setup, these events would simply accumulate in a source system and be pulled into the warehouse on a fixed schedule -- once an hour, once a day. Between runs, every dashboard and every downstream decision is working from data that is, on average, half of that interval out of date; right before the next run, it is a full interval out of date. For a platform that wants to react to what users are doing now -- surfacing trending songs, flagging a spike in failed logins, feeding a recommendation engine -- that lag is the problem: the data is correct, just always late.

Real-time (streaming) processing closes this gap. Instead of waiting for a scheduled run, each event is picked up and made queryable within seconds to minutes of happening, so dashboards and downstream models reflect what is happening now rather than what happened last interval. That said, doing this reliably, and not just quickly, is harder than it sounds once the platform has more than one producer and more than one consumer of the same data: events from different partitions can arrive out of order, any component in the chain can crash mid-stream, the same event must not be silently duplicated or dropped on retry, upstream fields will change shape over time, and the storage layer has to serve both a constant stream of small writes and large analytical reads without one starving the other. These are the distributed-systems problems that ordering guarantees, replication, exactly-once processing, and lakehouse table formats exist to solve, and they are exactly the concepts the rest of this chapter introduces: Apache Kafka for ordered, fault-tolerant ingestion, Spark Structured Streaming for continuous, exactly-once processing, and Apache Hudi for a storage layer that supports both frequent streaming writes and fast analytical reads.

\subsection{Event Streaming and Publish-Subscribe Messaging}

Event streaming treats data as an unbounded sequence of immutable events (e.g. ``user played a song'') rather than as static rows in a table. Producers append events to a log; one or more consumers read that log independently and at their own pace. Apache Kafka~\cite{kafka} implements this as a publish-subscribe (pub-sub) system: producers publish events to named \emph{topics}, and consumers subscribe to the topics they need, decoupling producers from consumers in time and in number.

\subsubsection{Partitioning and Replication}

A Kafka topic is split into \emph{partitions}, each an ordered, append-only log. Partitioning is Kafka's unit of parallelism: different partitions of the same topic can be produced to and consumed from concurrently, and ordering is only guaranteed within a single partition. Each partition can additionally be \emph{replicated} across multiple brokers (replication factor); one replica is the leader that serves reads and writes, and the others are followers that copy the leader's log. If a broker holding the leader replica fails, a follower on a surviving broker can take over, so the topic keeps serving as long as at least one replica of each partition survives.

\subsubsection{KRaft vs. ZooKeeper}

Earlier Kafka versions relied on Apache ZooKeeper, a separate coordination service, to elect partition leaders and store cluster metadata. KRaft (Kafka Raft) replaces that external dependency with a Raft-based consensus protocol run by the Kafka brokers themselves, so the brokers agree on cluster metadata directly instead of through a second distributed system. This removes an entire service (and its own quorum, ports, and failure modes) from the deployment, which matters for a small, resource-constrained cluster where every additional coordination process competes for the same limited memory and CPU.

\subsubsection{Consumer Groups and Offsets}

Kafka consumers read as part of a named \emph{consumer group}: Kafka assigns each partition of a subscribed topic to exactly one consumer within the group, so the group as a whole divides the work of consuming a topic. Each consumer tracks its \emph{offset} -- the position of the last record it has read -- per partition. Because a consumer's progress is expressed purely as an offset, a consumer (or its replacement) can resume from exactly where it left off after a restart, rather than re-reading the whole topic or losing its place, as long as the offset is durably recorded somewhere.

\subsection{Spark Structured Streaming and Micro-batching}

Apache Spark Structured Streaming~\cite{spark-structured-streaming} runs a streaming query as a series of small, discrete batch jobs (\emph{micro-batches}) on Spark's existing batch execution engine, rather than processing each event individually. On a fixed trigger interval, the engine reads whatever new data has arrived since the last batch (tracked via checkpointed offsets), applies the same DataFrame/SQL transformations used for batch queries, and writes the result. This design reuses Spark's mature batch engine (query planning, fault tolerance, exactly-once output via checkpointing) at the cost of end-to-end latency being bounded by the micro-batch interval rather than being event-at-a-time.

\subsubsection{Checkpointing and Exactly-once Output}

Structured Streaming persists two things to a durable checkpoint location after each micro-batch: which input offsets that batch covered, and the state needed to avoid reprocessing them. On restart, the engine reads the checkpoint to determine exactly which offsets have already been committed to the output, and resumes from there. Combined with an idempotent or transactional sink, this gives \emph{exactly-once} output semantics -- each input record affects the output exactly once, even across job restarts -- without the engine needing to coordinate with the source system beyond reading its offsets.

\subsection{Table Formats: Apache Hudi (MERGE\_ON\_READ vs. COPY\_ON\_WRITE)}

Apache Hudi~\cite{hudi} is a transactional table format for data lakes: it adds row-level insert/update/delete semantics, ACID commits, and time-travel to plain files on object storage (S3, MinIO, HDFS, ...). Hudi supports two table types that trade off write cost against read cost:

\begin{itemize}
    \item \textbf{COPY\_ON\_WRITE (COW).} Every write rewrites the affected data files as new versioned Parquet files. Reads are fast (plain Parquet, no merge step), but writes are more expensive because entire files are rewritten on update.
    \item \textbf{MERGE\_ON\_READ (MOR).} Writes are appended to a compact row-based delta log next to the existing base Parquet files; the two are merged only at read time (or during a background compaction job). Writes are cheap and low-latency, at the cost of extra work on read (or a periodic compaction step) to materialize the merged view.
\end{itemize}

MOR suits high-frequency streaming appends, where write latency matters most; COW suits tables that are rebuilt in full on a schedule and then read many times, where read performance matters most.

\subsubsection{Medallion Layering: Bronze and Gold}

The medallion architecture organizes a lakehouse into successive layers of increasing refinement, conventionally named bronze, silver, and gold. \emph{Bronze} tables hold raw or lightly-cleaned data close to its source form, written as-is by the ingestion layer. \emph{Gold} tables hold data that has been modeled and aggregated into the shape consumers (BI tools, dashboards) actually query, typically built by a transformation layer reading from bronze. This project uses the bronze and gold terms in that sense throughout: bronze is what Spark Structured Streaming writes directly from Kafka, and gold is what dbt builds on top of bronze.

\subsection{Star-Schema Dimensional Modeling}

Dimensional modeling organizes an analytical warehouse around a central \emph{fact table} of measurable events (here, one row per stream/play event) surrounded by \emph{dimension tables} of descriptive attributes (who, what, where, when) that the fact table references by key. This ``star'' shape -- one fact table joined to several dimension tables -- normalizes descriptive attributes into reusable dimensions while keeping the fact table narrow, and is the standard layout for BI and reporting workloads because it keeps common queries (aggregate the fact table, filter/group by a dimension) to a small, predictable number of joins.

\subsubsection{Slowly Changing Dimensions (Type 2)}

A dimension's attributes can change over time -- e.g. a user upgrading from a free to a paid subscription level -- and a \emph{Slowly Changing Dimension Type 2} keeps that history instead of overwriting it. Each change is stored as a new row with a validity window (an activation date and an expiration date), and the fact table joins to whichever dimension row was valid at the time of the event, rather than to a single current row. This preserves the ability to ask what a dimension's attributes were at any past point in time, not just what they are now.
